\section{Metodologia}\label{sec:metodologia}

Este capítulo descreve o ambiente experimental, o procedimento
de avaliação e os instrumentos utilizados para comparar o
controlador AIMD CUBIC com duas referências de base. A
organização segue seis eixos: o testbed físico e a carga de
trabalho adotada; os cenários de degradação de rede; os
controladores comparados; a instrumentação e as métricas
analisadas; o uso de inteligência artificial como ferramenta de
pesquisa; e o protocolo de execução das sessões experimentais.

\subsection{Testbed e Carga de Trabalho}

O servidor de \emph{streaming} executa o Sunshine sobre o
hardware descrito na Tabela~\ref{tab:testbed}. A escolha do
ambiente buscou aproximar-se de uma instalação típica de jogos
em nuvem \emph{self-hosted}: GPU dedicada da geração atual,
codificador NVENC, sistema operacional Linux, e conexão
cabeada ao roteador doméstico.

\begin{table}[!t]
  \caption{Especificações do testbed experimental.}
  \label{tab:testbed}
  \centering
  \begin{tabular}{ll}
    \toprule
    \textbf{Componente} & \textbf{Especificação} \\
    \midrule
    CPU (servidor) & AMD Ryzen 7 7700 (8 núcleos, 16 \emph{threads}) \\
    RAM (servidor) & 32~GB \\
    GPU (servidor) & NVIDIA GeForce RTX 4070 SUPER \\
    Codificador & NVENC \\
    Sistema operacional & Fedora Linux 43 \\
    \emph{Kernel} & Linux 6.19 \\
    Interface de rede (servidor) & Ethernet 1~Gbps \\
    Cliente & Smartphone Android, \emph{Moonlight} \\
    Conexão do cliente & Wi-Fi, mesma rede local \\
    \bottomrule
  \end{tabular}
\end{table}

A degradação de rede foi emulada no servidor via \texttt{tc}
com disciplina HTB, aplicada à interface de saída \texttt{eno1}.
Essa abordagem permite limitar a capacidade do enlace de forma
determinística e reproduzível, sem introduzir atrasos
artificiais que poderiam confundir o sinal de RTT utilizado
pelo controlador. Os \emph{scripts} de \emph{shaping} adotados estão disponíveis
no repositório do projeto.

A carga de trabalho utilizada em todas as sessões foi o jogo
\emph{Celeste}, escolhido por três motivos complementares. O
primeiro é o perfil de \emph{rendering} leve: o estilo
visual 2D mantém a taxa de bits demandada pelo codificador em
patamares relativamente estáveis, evitando que a variação do
canal seja confundida com variação do conteúdo. O segundo é a
sensibilidade do \emph{gameplay} à latência: a mecânica
baseia-se em saltos de precisão e janelas curtas de reação,
condição em que travamentos e atrasos são imediatamente
perceptíveis ao operador. O terceiro é a previsibilidade do
comportamento entre sessões, característica desejável para que
diferenças observadas entre controladores possam ser atribuídas
ao próprio controlador, e não ao conteúdo exibido. A escolha de
um único título constitui uma limitação reconhecida na
discussão dos resultados (Capítulo~\ref{sec:resultados}).

\subsection{Cenários de Avaliação}

Foram definidos seis cenários de degradação de rede,
sintetizados na Tabela~\ref{tab:scenarios}, escolhidos para
cobrir as principais formas em que um canal doméstico pode
degradar ao longo de uma sessão de jogo. Os cenários variam em
três dimensões complementares: a forma da degradação (queda
abrupta, em degraus, em janelas alternadas, ou contínua); a
severidade da constrição (de \SI{16}{\mega\bit\per\second}
até \SI{4}{\mega\bit\per\second}); e a duração total do
período sob shaping (de janelas curtas de \SI{30}{\second} a
constrições contínuas de \SI{160}{\second}). Fora das janelas
indicadas, o enlace opera sem restrição imposta pelo
\texttt{tc}, ou seja, na capacidade nativa de \SI{1}{\giga\bit\per\second}.

\begin{table}[!t]
  \caption{Cenários de degradação de rede utilizados na campanha
  experimental. Fora das janelas indicadas, o canal opera sem
  restrição.}
  \label{tab:scenarios}
  \centering
  \begin{tabular}{cll}
    \toprule
    \textbf{Cenário} & \textbf{Capacidade} & \textbf{Janela de shaping} \\
    \midrule
    A & 8~Mbps & 80--160~s \\
    B & 16 / 12 / 8~Mbps & 60--90 / 90--120 / 120--160~s \\
    C & 14 / 10~Mbps & 40--120 / 120--200~s \\
    D & 10~Mbps & 30--60 / 90--120 / 150--180 / 210--240~s \\
    E & 4~Mbps & 40--200~s \\
    F & 6~Mbps & 0--160~s \\
    \bottomrule
  \end{tabular}
\end{table}

O cenário~A representa uma queda abrupta e prolongada,
testando a capacidade do controlador de reagir a um corte
único e sustentado. O cenário~B aplica degraus sucessivos de
severidade crescente, exigindo do controlador adaptações
encadeadas em curto intervalo. O cenário~C aplica dois
patamares mais longos, simulando uma transição entre regimes
de banda. O cenário~D alterna entre quatro janelas curtas de
constrição e períodos de recuperação, testando o
comportamento do controlador em ciclos de degradação e
restauração. Os cenários~E e~F representam constrições
contínuas, severa e moderada respectivamente, em que o
controlador deve estabilizar a operação dentro do envelope
imposto pela rede.

\subsection{Controladores Comparados}

Três controladores foram avaliados, identificados como
referência fixa, AIMD de base e AIMD CUBIC. O controlador de
referência fixa corresponde ao comportamento original do
Sunshine, em que a taxa de bits é negociada uma única vez ao
início da sessão e mantida constante até o fim. O controlador
AIMD de base aplica corte multiplicativo com fator $\beta = 0{,}7$
ao detectar congestionamento e cresce em duas fases
(``$\times \gamma$ depois $+\Delta$''), sem o mecanismo de
\emph{fast convergence}; corresponde à versão do controlador
adaptativo que antecedeu a contribuição apresentada neste
trabalho. O controlador AIMD CUBIC corresponde à contribuição
descrita no Capítulo~\ref{sec:contribuicao}, com os parâmetros
sintetizados na Tabela~\ref{tab:params}.

\subsection{Instrumentação e Métricas}

A coleta de dados em \emph{runtime} é realizada pelo próprio
Sunshine, por meio de extensão de instrumentação introduzida em
contribuição prévia ao projeto. A cada amostra de RTT recebida
da biblioteca ENet, um registro é exportado para um arquivo
CSV contendo o instante de amostragem, o RTT, a variância do
RTT, o tempo de envio do quadro mais recente, o tamanho do
quadro codificado, a taxa de bits-alvo corrente, o estado do
controlador e indicadores de eventos relevantes (envio de
quadro IDR, requisição de IDR pelo cliente). A taxa de
amostragem é a própria taxa de quadros do codificador,
tipicamente próxima de \SI{60}{\hertz}.

A partir do CSV de cada sessão são derivadas as métricas
agregadas reportadas no Capítulo~\ref{sec:resultados}. As
métricas de latência incluem o pico de RTT, os percentis 50,
95 e 99 da distribuição, e a fração do tempo total em que o
RTT excede o limiar de \SI{100}{\milli\second} estabelecido
como referência para jogos de ação~\cite{baldovino2022}. As
métricas de throughput incluem a taxa de bits média sustentada
durante a janela de shaping e o uso do canal, definido como a
razão entre essa média e a capacidade imposta pela emulação.
As métricas de continuidade incluem a contagem de \emph{stalls},
definidos como intervalos consecutivos superiores a
\SI{100}{\milli\second} sem entrega de quadros ao cliente, o
tempo total acumulado em \emph{stall}, e a contagem de quadros
IDR emitidos durante a sessão. Esta última é utilizada como
indicador indireto de recuperação após perdas que escaparam aos
mecanismos de FEC.

Adicionalmente, ao final de cada sessão o operador registrou
três avaliações subjetivas em escala de 1 a 5, em que 1
corresponde à melhor avaliação: \emph{latência percebida},
\emph{qualidade visual} e \emph{estabilidade}. A avaliação foi
conduzida por um único operador (o autor do trabalho), o que
constitui uma limitação reconhecida na discussão dos
resultados.

\subsection{Uso de Inteligência Artificial como Ferramenta de Pesquisa}

Ferramentas de inteligência artificial baseadas em modelos de
linguagem de grande porte foram empregadas em dois pontos do
trabalho. O primeiro foi o apoio à implementação em C++ do
controlador e de sua integração ao Sunshine, incluindo o
ajuste de testes em \emph{gtest} e a conformidade com as
convenções de estilo do projeto \emph{upstream}. Todo o código
produzido com esse apoio foi revisado, compilado e testado pelo
pesquisador antes de qualquer submissão ao repositório
\emph{upstream}, e nenhuma alteração foi incorporada sem
verificação.

O segundo foi a estruturação de uma base de conhecimento
local, em formato Markdown compatível com o \emph{Obsidian},
utilizada como apoio à navegação do material bibliográfico, da
documentação interna do Sunshine, e do registro das iterações
de \emph{design} do controlador. A base é composta por
análises individuais (de \emph{papers} lidos, de subsistemas
do Sunshine e de Linux, e de cada iteração de projeto), com
referências cruzadas via \emph{wikilinks} que permitem
visualização em grafo. A síntese inicial de cada arquivo foi
auxiliada por IA a partir da leitura primária do material pelo
pesquisador, com edições e validações subsequentes realizadas
manualmente.

O uso de IA neste trabalho restringiu-se ao apoio à
implementação de código e à estruturação da base de
conhecimento de pesquisa, ambos sob revisão e validação do
pesquisador. As decisões de projeto, a análise dos resultados
experimentais e as conclusões apresentadas são de autoria
própria.

\subsection{Protocolo de Execução}

A campanha experimental consistiu em uma execução por par
(controlador, cenário) para o controlador AIMD CUBIC,
totalizando seis execuções novas. As colunas correspondentes
ao controlador de referência fixa e ao AIMD de base foram
reaproveitadas de campanha anterior, conduzida sobre o mesmo
testbed e com a mesma configuração do servidor, em razão da
baixa variância observada entre execuções repetidas naquela
campanha. As implicações dessa escolha são discutidas no
Capítulo~\ref{sec:resultados}.

Cada execução seguiu o mesmo procedimento. O servidor foi
iniciado com a configuração apropriada ao controlador em
avaliação. O cliente conectou-se ao servidor e iniciou o jogo
\emph{Celeste}, com o operador conduzindo o personagem por uma
sequência de fases pré-determinada, de duração total de
aproximadamente quatro minutos, suficiente para cobrir a janela
de shaping mais períodos de aquecimento e encerramento. O \emph{shaping} foi acionado no instante zero
da contagem do cenário, com a duração definida na
Tabela~\ref{tab:scenarios}. Ao final da sessão, o operador
registrou as três avaliações subjetivas descritas na seção
anterior. O arquivo CSV resultante foi arquivado em
\texttt{testbed/results/cubic/} com identificador associado ao
cenário e ao instante de execução.
